
\clearpage
\section{Realizaciones elementales del Modelo Estándar}
\begingroup
\footnotesize
\setlength{\tabcolsep}{1.9pt}
\begin{longtable}{L{2.03cm}L{2.03cm}L{1.87cm}L{2.34cm}L{6.55cm}}
\toprule
\rowcolor{HMTNavy}\HMTtablehead{Partícula} & \HMTtablehead{Carga} & \HMTtablehead{Espín} & \HMTtablehead{Estadística} & \HMTtablehead{Realización HMT--MD} \\
\midrule
\endfirsthead
\multicolumn{5}{l}{\sffamily\scriptsize\color{HMTGray}Continuación de la tabla}\\[-1pt]
\toprule
\rowcolor{HMTNavy}\HMTtablehead{Partícula} & \HMTtablehead{Carga} & \HMTtablehead{Espín} & \HMTtablehead{Estadística} & \HMTtablehead{Realización HMT--MD} \\
\midrule
\endhead
\midrule
\endfoot
\bottomrule
\endlastfoot
\(e^-\) & \(-1\) & \(1/\allowbreak{}2\) & fermiónica & punto fijo \((5,5)\), vacancia orientada \\
\rowcolor{HMTLight}\(\mu^-\) & \(-1\) & \(1/\allowbreak{}2\) & fermiónica & bloque leptónico de rango 8 \\
\(\tau^-\) & \(-1\) & \(1/\allowbreak{}2\) & fermiónica & bloque leptónico de rango 16 \\
\rowcolor{HMTLight}\(\nu_e\) & \(0\) & \(1/\allowbreak{}2\) & fermiónica & vector de interacción en soporte neutro 20 \\
\(\nu_\mu\) & \(0\) & \(1/\allowbreak{}2\) & fermiónica & vector de interacción en soporte neutro 20 \\
\rowcolor{HMTLight}\(\nu_\tau\) & \(0\) & \(1/\allowbreak{}2\) & fermiónica & vector de interacción en soporte neutro 20 \\
\(u\) & \(+2/\allowbreak{}3\) & \(1/\allowbreak{}2\) & fermiónica & rama up, generación 1, triplete de color \\
\rowcolor{HMTLight}\(d\) & \(-1/\allowbreak{}3\) & \(1/\allowbreak{}2\) & fermiónica & rama down, generación 1, triplete de color \\
\(c\) & \(+2/\allowbreak{}3\) & \(1/\allowbreak{}2\) & fermiónica & rama up, generación 2, triplete de color \\
\rowcolor{HMTLight}\(s\) & \(-1/\allowbreak{}3\) & \(1/\allowbreak{}2\) & fermiónica & rama down, generación 2, triplete de color \\
\(t\) & \(+2/\allowbreak{}3\) & \(1/\allowbreak{}2\) & fermiónica & rama up, generación 3, triplete de color \\
\rowcolor{HMTLight}\(b\) & \(-1/\allowbreak{}3\) & \(1/\allowbreak{}2\) & fermiónica & rama down, generación 3, triplete de color \\
\(\gamma\) & \(0\) & \(1\) & bosónica & sección nula, dos polarizaciones \\
\rowcolor{HMTLight}\(g^a\) & \(0\) & \(1\) & bosónica & adjunta de color, ocho generadores \\
\(W^\pm\) & \(\pm1\) & \(1\) & bosónica & bloque vectorial cargado conjugado \\
\rowcolor{HMTLight}\(Z^0\) & \(0\) & \(1\) & bosónica & bloque vectorial neutral \\
\(H\) & \(0\) & \(0\) & bosónica & sección escalar persistente \\
\end{longtable}
\endgroup
Las antipartículas se obtienen por \(C_M\). Los antiquarks pertenecen al triplete conjugado; \(W^-\) es el conjugado de \(W^+\); fotón, \(Z\), Higgs y gluones poseen las conjugaciones propias de sus representaciones.

El soporte B1 exacto de las clases elementales tiene rangos

\[
 1,\quad8,\quad16,\quad
 3\times20,\quad
 3\times16,\quad
 3\times20.
\]
Por tanto, el número de incidencias clase--ruta es

\[
 1+8+16+60+48+60=193.
\]
Las incidencias no son 193 rutas distintas: los tres sabores neutrínicos comparten el soporte neutro de rango veinte; los sabores up reutilizan un soporte de rango dieciséis y los sabores down, uno de rango veinte. El proyector de sabor debe refinar estos soportes. En la unión quark de 36 celdas, el color queda exactamente

\[
 12_R+12_G+12_B.
\]
Esta igualdad certifica el soporte cromático global, pero no autoriza todavía a elegir una celda por sabor.

\Needspace{7\baselineskip}
\section{Descomposición del sector leptónico por generación, carga, espín y carácter de masa}
\Needspace{7\baselineskip}
\subsection{Electrón: centro, vacancias y orientación}
El electrón es excepcional porque su fibra contiene el único punto fijo \((5,5)\). Su reducción a masa escalar no requiere romper una órbita. La ruta activa es

\[
 e_{\rm act}=(2,2,5,-4,-3,-2)^2,
\]
mientras que la envolvente histórica

\[
 e_{\rm env}=(4,3,5,-4,-3,-5)^2
\]
conserva signo y máximo absoluto de las dos hojas. La envolvente retiene la palabra

\[
 \tau_e=+++---+++---,
\]
pero no sustituye la historia activa.

Cuatro datos que deben mantenerse distintos son:

\[
 \text{firma cruda}=(24,0,12,0,-12),
\]
\[
 \text{firma par CT--108}=(-12,-4,12,-6,12),
\]
\[
 \text{vector de acontecimientos}=e_{\rm act},
\qquad
 \text{origen afín}=v_e=0.
\]
El selector basal \((-22,+15)\) posee dos lecturas internas complementarias. La descomposición regional da

\[
 -22=-(27-5),\qquad +15=5\cdot3=\binom62,
\]
y el marco de vacancias

\[
 V_e=
 \begin{pmatrix}
 21&22\\
 6&7
 \end{pmatrix}
\]
satisface

\[
 \det V_e=21\cdot7-22\cdot6=15.
\]
El \(22\) mide el déficit orientado de la extensión; el \(15\), la retención de área que sobrevive al cierre. Así, la partícula electrónica aparece como persistencia central de una vacancia orientada, no como una masa insertada en la celda central.

La ecuación completa es

\[
 \mathfrak e_e^\beta=
 \frac{\sqrt3}{4}
 \left[
 \exp\!\left(\frac{\varphi}{\pi^2}\right)
 -22\alpha_{\rm HMT}^{3}
 \right]
 (1+15\Delta_4)
 R_{\rm act}^{80-54\alpha_{\rm HMT}+6\Delta_4}.
\]
En la carta electrónica,

\[
 \mathfrak e_e^\beta
 =0.5109989506900008086082571648\ldots\ {\rm MeV}.
\]
\Needspace{7\baselineskip}
\subsection{Realización del muón como sección leptónica de 2.ª generación: fibra, firma y carácter}
El muón reside en una multisección leptónica de rango ocho. Su huella estructural recuperada comprende las celdas

\[
 \{21,23,25,39,43,57,59,61\},
\]
los conteos de giro/cambio \(22/25\), el eje norte--sur,

\[
 (N,E,S,O)=(62,25,7,14),
\qquad
 (h_{369},r_5)=(57,9),
\]
y la coordenada central

\[
 (n_A,s_C)=(42,-3).
\]
El objeto propio es el operador comprimido

\[
 \widehat M_\mu=P_\mu\widehat M P_\mu.
\]
Su reducción a un número exige probar que el proyector de sabor y generación conmuta con el operador, o que el estado preparado tiene varianza nula. Dada la firma enriquecida

\[
 (42,-3,8,0,2),
\]
la evaluación del carácter es

\[
 105.658360294435829\ldots\ {\rm MeV}.
\]
Esta cifra es un control exacto de la evaluación de esa firma; la derivación física completa debe reproducir la firma desde el registro genealógico seleccionado por el acoplamiento. La anchura muónica se obtiene del operador temporal de decaimiento y no de \(\widehat M_\mu\).

\Needspace{7\baselineskip}
\subsection{Realización del tau como sección leptónica de 3.ª generación: fibra, firma y carácter}
El tau ocupa una multisección leptónica de rango dieciséis. Su huella comprende

\[
 17/24,\qquad \text{eje este--oeste},\qquad
 (N,E,S,O)=(29,55,7,17),
\]
\[
 (h_{369},r_5)=(59,6),\qquad
 (n_A,s_C)=(64,0).
\]
La firma enriquecida

\[
 (64,0,25,-1,6)
\]
produce

\[
 1776.8609176314323\ldots\ {\rm MeV}.
\]
Como en el muón, la estructura firme es el bloque operatorial; la masa escalar se obtiene cuando el acoplamiento tau selecciona un bloque escalar o un polo aislado.

\Needspace{7\baselineskip}
\subsection{Fibra neutrínica común y separación de sabores mediante el lector PMNS}
Los tres sabores neutrínicos comparten el espacio

\[
 \mathcal H_\nu=
 \bigoplus_{j=1}^{4}\mathcal H_{\nu,G_j},
\qquad
 \dim\mathcal H_\nu=2+4+6+8=20.
\]
Los vectores

\[
 f_e,\quad f_\mu,\quad f_\tau
\]
constituyen una base de interacción; no son tres autovectores de masa. La profundidad \(G_j\) tampoco es un cuarto sabor estéril: una interpretación estéril requiere un proyector físico adicional.

El extractor neutro se define por

\[
 T_\nu(m)=\sum_{t\in E_m}
 \omega(t)\varepsilon(t)\theta(t),
\qquad
 e_m^\nu=T_\nu(m)-T_\nu(m+6),
\]
\[
 \tau_m^\nu=\operatorname{sgn}(e_m^\nu).
\]
De él se obtienen

\[
 b_{120}^\nu
 =\sum_{a=0}^{3}
 \operatorname{sgn}\!\left(\sum_{m\in I_a}e_m^\nu\right),
\]
\[
 b_\zeta^\nu=\sum_m\tau_m^\nu\tau_{m+3}^\nu.
\]
La firma neutrínica posee la forma

\[
 v_\nu^{(i)}
 =(0,0,k_{\Delta_4}^{(i)},b_{120}^{\nu_i},b_\zeta^{\nu_i}).
\]
La inversión bidireccional cambia \(b_{120}\) y conserva \(b_\zeta\); esta asimetría produce internamente las dos orientaciones de jerarquía.

\Needspace{7\baselineskip}
\subsection{Naturaleza fermiónica del neutrino}
La neutralidad no determina la estadística. El sector neutrínico vive en la completación exterior y porta la holonomía espinorial:

\[
 U_\nu(2\pi)=-I,\qquad U_\nu(4\pi)=I,
\]
\[
 \{a_{\nu,i},a_{\nu,j}^\dagger\}=\delta_{ij}.
\]
Por tanto, cada modo neutrínico físico es fermiónico. La pregunta Dirac/Majorana es posterior: decide si la conjugación intercambia dos secciones distintas o identifica una sección con su conjugada dentro de una representación real. Ninguna de las dos posibilidades convierte el neutrino en bosón.

\Needspace{7\baselineskip}
\subsection{Operador de masa y PMNS}
Sea \(P_\nu\) el proyector de rango tres sobre el subespacio de modos másicos. Primero se diagonaliza

\[
 P_\nu\widehat M_\nu P_\nu
 =\sum_{i=1}^{3}m_i\,|n_i\rangle\langle n_i|.
\]
Después se construye el transporte entre bases:

\[
 (f_e,f_\mu,f_\tau)
 =(n_1,n_2,n_3)U_{\rm PMNS}^\dagger.
\]
Un núcleo de rango completo que lee hoja, orientación, acarreo, vacancia y holonomía puede escribirse

\[
 G_{ab}=
 \frac1{\sqrt{|L_a||N_b|}}
 \sum_{c\in L_a}\sum_{d\in N_b}K_{\rm reg}(c,d),
\]
\[
 U_{\rm PMNS}=\operatorname{polar}(G),
\qquad \det G\ne0.
\]
La fase interna de rango uno produce los ángulos

\[
 (19.6807^\circ,56.4970^\circ,29.6224^\circ)
\]
y, por ello, no puede ser el núcleo PMNS físico completo. El rango completo es una exigencia matemática, no un retoque numérico.

Los observables

\[
 m_1,m_2,m_3,\qquad
 \Delta m_{ij}^2,\qquad
 m_\beta,\qquad
 m_{\beta\beta},\qquad
 \sum_i m_i
\]
son mapas diferentes del mismo espectro. No deben colapsarse en una sola “masa del neutrino”.

\Needspace{7\baselineskip}
\section{Quarks, color, generaciones y mezcla CKM}
\Needspace{7\baselineskip}
\subsection{Posición estructural de los \(6\) sabores de quark}
Cada quark se define por

\[
 \mathfrak q_X=
 (\text{rama de carga},\text{generación},
 \mathbb Z_3^{\rm color},\text{multisección},
 \text{registro genealógico},\text{hoja}).
\]
La posición estructural de los seis sabores queda resumida por sus huellas anteriores a toda carta de masa:

\Needspace{7\baselineskip}
\begingroup
\fontsize{7.2}{8.2}\selectfont
\setlength{\tabcolsep}{1.9pt}
\begin{longtable}{L{0.94cm}L{2.81cm}L{2.03cm}L{1.87cm}L{3.28cm}L{2.18cm}L{1.87cm}}
\toprule
\rowcolor{HMTNavy}\HMTtablehead{Sabor} & \HMTtablehead{Rama/\allowbreak{}generación} & \HMTtablehead{Giro/\allowbreak{}cambio} & \HMTtablehead{Eje} & \HMTtablehead{\((N,E,S,O)\)} & \HMTtablehead{\((h_{369},r_5)\)} & \HMTtablehead{\((n_A,s_C)\)} \\
\midrule
\endfirsthead
\multicolumn{7}{l}{\sffamily\scriptsize\color{HMTGray}Continuación de la tabla}\\[-1pt]
\toprule
\rowcolor{HMTNavy}\HMTtablehead{Sabor} & \HMTtablehead{Rama/\allowbreak{}generación} & \HMTtablehead{Giro/\allowbreak{}cambio} & \HMTtablehead{Eje} & \HMTtablehead{\((N,E,S,O)\)} & \HMTtablehead{\((h_{369},r_5)\)} & \HMTtablehead{\((n_A,s_C)\)} \\
\midrule
\endhead
\midrule
\endfoot
\bottomrule
\endlastfoot
\(u\) & positiva/\allowbreak{}1 & \(21/\allowbreak{}30\) & N--S & \((34,31,31,12)\) & \((51,13)\) & \((11,7)\) \\
\rowcolor{HMTLight}\(d\) & negativa/\allowbreak{}1 & \(20/\allowbreak{}32\) & E--O & \((27,35,18,28)\) & \((47,17)\) & \((17,8)\) \\
\(c\) & positiva/\allowbreak{}2 & \(21/\allowbreak{}29\) & E--O & \((20,38,24,26)\) & \((49,16)\) & \((61,8)\) \\
\rowcolor{HMTLight}\(s\) & negativa/\allowbreak{}2 & \(17/\allowbreak{}27\) & N--S & \((50,32,12,14)\) & \((57,11)\) & \((41,-3)\) \\
\(t\) & positiva/\allowbreak{}3 & \(21/\allowbreak{}26\) & N--S & \((46,27,16,19)\) & \((57,9)\) & \((100,-1)\) \\
\rowcolor{HMTLight}\(b\) & negativa/\allowbreak{}3 & \(21/\allowbreak{}30\) & E--O & \((12,63,9,24)\) & \((47,7)\) & \((71,-5)\) \\
\end{longtable}
\endgroup
Estas posiciones no son seis números elegidos: son seis clases de representación que actúan sobre tripletes de color. La ruta física puede ser una órbita o un bloque de rutas; no tiene que ser un punto.

\Needspace{7\baselineskip}
\subsection{Invarianza escalar del carácter de masa bajo la acción de color}
Sea \(P_q\) el proyector de sabor y generación y \(P_{\rm RGB}\) el proyector a la órbita de color. La masa quark es

\[
 \widehat M_q=
 P_qP_{\rm RGB}\widehat M P_{\rm RGB}P_q.
\]
Si

\[
 [\widehat M_q,U(g)]=0,\qquad g\in SU(3),
\]
entonces

\[
 \widehat M_q=m_qI_3
\]
sobre el triplete irreducible. Esta igualdad explica por qué el sabor porta una masa común aunque sus tres colores ocupen posiciones distintas en el atlas.

\Needspace{7\baselineskip}
\subsection{Carta de esquema y escala}
La coordenada HMT interna debe transportarse a una definición renormalizada:

\[
 m_q^{\mathcal S}(\mu)
 =Z_m^{\mathcal S}(\mu,\mu_0)\,
 \mathcal R_q(\widehat M_q),
\]
con ley de composición

\[
 Z_m(\mu_2,\mu_0)
 =Z_m(\mu_2,\mu_1)Z_m(\mu_1,\mu_0).
\]
Para \(u,d,s\) la comparación natural usa \(\overline{\rm MS}\) a \(2\,{\rm GeV}\); para \(c,b\), la escala de la propia masa renormalizada; para \(t\), debe declararse si la carta es de polo, MSR o cinemática. Sin este transporte, una diferencia numérica no es todavía un residual de la ley.

\Needspace{7\baselineskip}
\subsection{Evaluaciones de las firmas declaradas}
Las coordenadas centrales históricas producen:

\Needspace{7\baselineskip}
\begingroup
\footnotesize
\setlength{\tabcolsep}{2.1pt}
\begin{longtable}{L{4.68cm}L{10.61cm}}
\toprule
\rowcolor{HMTNavy}\HMTtablehead{Sabor} & \HMTtablehead{Evaluación} \\
\midrule
\endfirsthead
\multicolumn{2}{l}{\sffamily\scriptsize\color{HMTGray}Continuación de la tabla}\\[-1pt]
\toprule
\rowcolor{HMTNavy}\HMTtablehead{Sabor} & \HMTtablehead{Evaluación} \\
\midrule
\endhead
\midrule
\endfoot
\bottomrule
\endlastfoot
\(u\) & \(2.165924\allowbreak{}536173\allowbreak{}907\ldots\ {\rm MeV}\) \\
\rowcolor{HMTLight}\(d\) & \(4.679550\allowbreak{}162948\allowbreak{}874\ldots\ {\rm MeV}\) \\
\(c\) & \(1.270500\allowbreak{}838641\allowbreak{}911\ldots\ {\rm GeV}\) \\
\rowcolor{HMTLight}\(s\) & \(92.940093\allowbreak{}907845\allowbreak{}64\ldots\ {\rm MeV}\) \\
\(t\) & \(172.597479\allowbreak{}544019\allowbreak{}1\ldots\ {\rm GeV}\) \\
\rowcolor{HMTLight}\(b\) & \(4.190119\allowbreak{}763299\allowbreak{}790\ldots\ {\rm GeV}\) \\
\end{longtable}
\endgroup
Estas evaluaciones comprueban el carácter una vez dada la firma. La deducción completa exige obtener cada firma, su refinamiento de torre y su carta de esquema a partir del proyector de sabor y del registro genealógico correspondiente.

\Needspace{7\baselineskip}
\subsection{Determinación de los ángulos CKM desde la reflexión racional de las rutas quark}
Sean

\[
 A=1000\alpha_{\rm HMT},\qquad
 h=\frac{C^*}{6},\qquad
 \Delta=\frac{180}{\pi}\Delta_4.
\]
La gramática angular da

\[
 \begin{pmatrix}
 \theta_{12}\\
 \theta_{23}\\
 \theta_{13}
 \end{pmatrix}
 =
 \begin{pmatrix}
 2&-5&28\\
 0&7&-25/2\\
 1&-20&-2/3
 \end{pmatrix}
 \begin{pmatrix}
 A\\h\\\Delta
 \end{pmatrix},
\qquad
 \delta_{\rm CP}=9A.
\]
El falsador lineal exacto es

\[
 3857\,\delta_{\rm CP}
 =9\bigl(
 1528\theta_{12}
 +3380\theta_{23}
 +801\theta_{13}
 \bigr).
\]
A partir de los tres ángulos y de \(\delta_{\rm CP}\) se construye la matriz unitaria \(V_{\rm CKM}\). Su función en la ley de masas no es producir autovalores, sino transportar los operadores down a la base up:

\[
 B_d^{(u)}=V_{\rm CKM}B_dV_{\rm CKM}^*.
\]
\Needspace{7\baselineskip}
\subsection{Invariante de Jarlskog}
El invariante de violación CP es

\[
 J_{\rm CKM}
 =c_{12}c_{23}c_{13}^2
 s_{12}s_{23}s_{13}\sin\delta_{\rm CP},
\]
y la evaluación interna es

\[
 J_{\rm CKM}
 =3.1189723326632947\times10^{-5}.
\]
Equivalentemente, la compatibilidad entre espectros y mezcla puede leerse en el determinante del conmutador de los operadores hermíticos up/down. Un valor no nulo de \(J_{\rm CKM}\) registra orientación física no eliminable; no es un corrector de masa añadido a cada quark.

\Needspace{7\baselineskip}
\section{Sectores de calibre y escalar}
\Needspace{7\baselineskip}
\subsection{Carta del sector nulo: rutas de masa \(0\), representaciones de calibre y persistencia}
Sea

\[
 \mathcal A_{\rm sp}=\mathbb R[R]/(R^2-1),
 \qquad
 P_\pm=\frac{1\pm R}{2}.
\]
Para

\[
 x=r_+P_++r_-P_-,
\]
la norma partida es

\[
 N_{\rm sp}(x)=r_+r_-.
\]
La sección nula es una dirección de norma cero. Su masa es exactamente

\[
 m=0,
\]
sin pasar por un límite de caracteres positivos y sin asignar \(\kappa=0\). La nulidad pertenece al tipo de la sección.

\Needspace{7\baselineskip}
\subsection{Realización del fotón en el sector nulo electromagnético: helicidad, calibre y carácter}
El fotón es la realización vectorial transversal de una sección de calibre nula. El proyector físico debe dejar exactamente dos polarizaciones:

\[
 \operatorname{rg}P_\gamma^{\rm trans}=2,
\qquad
 P_\gamma^{\rm trans}k=0.
\]
La masa de reposo es

\[
 m_\gamma=0.
\]
La propagación bidireccional puede portar una anisotropía orientada

\[
 \epsilon=q_-^{90}-q_+^{90},
\qquad
 c_\pm=\frac{c}{1\pm\epsilon},
\]
mientras el promedio armónico conserva

\[
 \frac{2}{1/c_++1/c_-}=c.
\]
Así, la ida y la vuelta pueden diferir sin alterar el invariante two--way. Este sesgo no es una masa de fotón.

Una realización Proca/Stückelberg distinta aparece si un brecha espectral del registro genealógico satisface

\[
 \omega^2=k^2+m_T^2
\]
y existe un entrelazador dimensional

\[
 m_T\longmapsto m_\gamma.
\]
La demora inducida obedecería

\[
 \frac{\Delta T}{T}
 \simeq\frac12
 \left(\frac{m_\gamma c^2}{\hbar\omega}\right)^2.
\]
El sector Maxwell nulo y la realización Proca masiva son dos ramas distintas; una no invalida a la otra.

\Needspace{7\baselineskip}
\subsection{Realización cromodinámica del sector nulo: octete de gluones y representación adjunta de \(SU(3)\)}
El trit residual genera el espacio de color, cuya parte sin traza realiza

\[
 \mathfrak{sl}_3\longrightarrow\mathfrak{su}_3.
\]
El gluón ocupa la representación adjunta de dimensión ocho. Mientras la simetría de color permanezca exacta,

\[
 m_g=0.
\]
Confinamiento, apantallamiento y brecha espectral colectivo son propiedades del espectro interactivo; no constituyen una masa de polo elemental del gluón. Una masa efectiva debe indicar observable, calibre, régimen y operador que la define.

\Needspace{7\baselineskip}
\subsection{Realización de los bosones \(W^\pm\): carga, polarización y carácter electrodébil}
El sector \(W\) es un bloque vectorial cargado con conjugación

\[
 C_M:W^+\longleftrightarrow W^-.
\]
Su huella histórica es

\[
 15/37,\qquad \text{eje este--oeste},\qquad
 (n_A,s_C)=(94,-1),
\]
y la firma enriquecida declarada es

\[
 (94,-1,0,-2,4).
\]
La evaluación correspondiente es

\[
 80.37896490970455\ldots\ {\rm GeV}.
\]
La realización completa es el polo del resolvente del bloque \(P_W\widehat M P_W\); su parte imaginaria determina la anchura. La degeneración de \(W^+\) y \(W^-\) procede de la conjugación del operador.

\Needspace{7\baselineskip}
\subsection{Realización del bosón \(Z^0\): mezcla neutra, polarización y carácter electrodébil}
El \(Z^0\) es el bloque vectorial neutral del acoplamiento electrodébil. Su huella histórica es

\[
 20/23,\qquad \text{eje este--oeste},\qquad
 (n_A,s_C)=(95,-1),
\]
y la firma enriquecida

\[
 (95,-1,-11,0,-4).
\]
La evaluación del carácter da

\[
 91.18753344431341\ldots\ {\rm GeV}.
\]
El proyector neutral debe construirse antes de la carta de polo; masa y anchura son dos coordenadas del mismo polo, no una única cifra.

\Needspace{7\baselineskip}
\subsection{Realización del sector de Higgs: modo escalar, acoplamiento y carácter de masa}
El Higgs es una sección escalar persistente del bloque de calibre con paridad CP declarada. Su huella es

\[
 24/29,\qquad \text{eje norte--sur},\qquad
 (n_A,s_C)=(97,9).
\]
La evaluación condicionada del carácter produce

\[
 125.2924660425361\ldots\ {\rm GeV}.
\]
La definición HMT no identifica “escalar” con “masa”: el espín cero fija la representación rotacional; la masa procede del espectro del bloque escalar acoplado. La anchura y los canales de decaimiento requieren el operador temporal.

\Needspace{7\baselineskip}
\subsection{El momento magnético anómalo del muón}
El funcional interno

\[
 a_\mu=
 \frac{\alpha_{\rm HMT}}{2\pi}
 +\frac{\alpha_{\rm HMT}(\varphi-1)}{1000}
 +\frac{\alpha_{\rm HMT}^2}{1000(54+1)}
\]
da

\[
 a_\mu=0.001165920713008014\ldots,
\]
y

\[
 g_\mu=2(1+a_\mu)
 =2.002331841426016028\ldots.
\]
Los tres términos corresponden, respectivamente, al giro basal, al retorno orientado en la completación \(1000\) y a la corrección cuadrática mediada por el defecto \(54\). Este observable comprueba la estructura interna del sector muónico, pero no selecciona por sí solo su ruta de masa.

\Needspace{7\baselineskip}
\section{Construcción de partículas compuestas mediante composición no aditiva de rutas y enlaces}
\Needspace{7\baselineskip}
\subsection{Operador de composición}
La composición de rutas no es estrictamente aditiva porque frontera, orden temporal y acarreo interactúan. Para dos registros,

\[
 \mathcal L_1\star\mathcal L_2
 =\mathcal L_1+\mathcal L_2
 +\mathfrak b(\mathcal L_1,\mathcal L_2),
\]
donde \(\mathfrak b\) es el cociclo de ligadura. La asociatividad exige

\[
 \mathfrak b(L_1,L_2)
 +\mathfrak b(L_1\star L_2,L_3)
 =
 \mathfrak b(L_2,L_3)
 +\mathfrak b(L_1,L_2\star L_3).
\]
El lector de masa actúa después de esta composición:

\[
 (\Gamma_i,\mathcal T,\partial)
 \xrightarrow{\mathfrak C_{12}}
 \mathcal L_{\rm comp}
 \xrightarrow{L_{\rm tick}}
 \sigma_{\rm comp}
 \xrightarrow{\chi_{\rm full}}
 \widehat M_{\rm comp}.
\]
\Needspace{7\baselineskip}
\subsection{Realización mesónica mediante composición quark–antiquark y ley cocíclica del enlace}
El dominio mesónico contiene 432 expresiones orientadas

\[
 M(q,q')=q\otimes(q')^\vee
\]
con color compatible. El proyector físico impone singlete de color y fija

\[
 (J^{PC},I,I_3,S,C,B,\ldots),
\]
junto con el árbol temporal de enlace. La masa de un mesón no es \(m_q+m_{\bar q}\); es un autovalor del registro compuesto.

Las seis evaluaciones históricas de control incluyen

\[
 m_{\pi^0}=134.9767243278721\ldots\ {\rm MeV},
\]
\[
 m_{\pi^\pm}=139.5704851715722\ldots\ {\rm MeV},
\]
\[
 m_{K^\pm}=493.6770856495306\ldots\ {\rm MeV},
\]
\[
 m_{K^0}=497.6111864977505\ldots\ {\rm MeV}.
\]
Estas cifras son evaluaciones exactas de registros declarados. El cierre genealógico estado por estado requiere regenerar cada registro mediante \(\mathfrak C_{12}\), incluida su frontera de ligadura.

\Needspace{7\baselineskip}
\subsection{Realización bariónica mediante composición trilineal de rutas quark y neutralización de color}
El dominio bariónico contiene 1728 expresiones RGB. El proyector físico debe imponer simultáneamente:

\begin{enumerate}
\item singlete de color;
\item estadística fermiónica;
\item sabor e isospín;
\item espín y simetrización;
\item árbol triple con acarreo y memoria.
\end{enumerate}
Para protón y neutrón, las firmas históricas

\[
 \sigma_p=(59,0,9,1,0),\qquad
 \sigma_n=(59,1,-34,0,1)
\]
dan

\[
 m_p=938.2717515469849\ldots\ {\rm MeV},
\]
\[
 m_n=939.5658104725070\ldots\ {\rm MeV}.
\]
La precisión aparente no sustituye la derivación del registro genealógico de ligadura. El signo y la magnitud de

\[
 m_n-m_p
\]
deben salir de orientación, isospín, vacancias y frontera de enlace antes de abrir la comparación externa.

\Needspace{7\baselineskip}
\subsection{Estados excitados y resonancias hadrónicas}
Una identidad compuesta puede poseer varios polos. El estado físico queda individualizado por

\[
 (\text{constituyentes},\mathcal T,\partial,
 J^{PC},I,\text{radial},\text{orbital},\text{canal}).
\]
La misma composición de sabores no basta para distinguir excitaciones. Cada polo

\[
 z_k=M_k-\frac{i}{2}\Gamma_k
\]
debe proceder de un autovector del operador compuesto y de un canal de decaimiento. Esto extiende la ley desde unas pocas masas estables a todo el espectro de mesones y bariones sin convertir el catálogo externo en generador.

\Needspace{7\baselineskip}
\subsection{Inventario exhaustivo de anchuras y tasas de transición de las realizaciones del atlas}
Sea \(P\) el subespacio preparado y \(Q=I-P\). El Hamiltoniano efectivo es

\[
 H_{\rm eff}(z)
 =PHP+PHQ(z-QHQ)^{-1}QHP.
\]
Los polos de

\[
 \det(z-H_{\rm eff}(z))=0
\]
determinan masas y anchuras. El carácter de masa aporta la parte real basal del operador; el término de autoenergía de \(Q\) aporta desplazamiento y decaimiento. Por ello, las 384 anchuras del inventario requieren un lector de canales adicional al lector de las 471 masas.

\Needspace{7\baselineskip}
\section{Sectores topológicos y partículas hipotéticas}
\Needspace{7\baselineskip}
\subsection{Sector topológico de Fibonacci: anillo basado, fusión y representación de trenzas}
La incidencia areal--volumétrica

\[
 F_{\rm av}=
 \begin{pmatrix}
 0&1\\1&1
 \end{pmatrix}
\]
realiza el anillo

\[
 \tau^2=1+\tau,\qquad
 \operatorname{FPdim}\tau=\varphi.
\]
La salida primaria es una regla de fusión y una dimensión cuántica. Una partícula anyónica de Fibonacci exige además matrices \(F\) y \(R\), pentágono, hexágono, quiralidad, Hamiltoniano con brecha espectral y un mapa desde el estado TPK al espacio físico. Sólo entonces puede definirse

\[
 m_{\rm brecha espectral}=\frac{\Delta E}{c_{\rm eff}^2}.
\]
\Needspace{7\baselineskip}
\subsection{Sector de Ising–Majorana: objetos simples, reglas de fusión y representación de trenzas}
La categoría de Ising satisface

\[
 \psi^2=1,\qquad
 \psi\sigma=\sigma,\qquad
 \sigma^2=1+\psi,\qquad
 d_\sigma=\sqrt2.
\]
Deben distinguirse cuatro niveles:

\begin{enumerate}
\item la involución de rutas \(C_M=-\operatorname{rev}\);
\item la paridad de la completación CAR;
\item un espinor relativista autorconjugado;
\item un modo cero topológico de Ising.
\end{enumerate}
Un fermión Majorana relativista requiere una representación espinorial real, una dinámica y un término de masa. Un modo cero de Majorana se caracteriza por brecha espectral, splitting y protección; no posee una masa universal deducible sólo del anillo de fusión. La identidad de una ruta con su conjugada es condición necesaria, pero no suficiente, para la realización Majorana.

\Needspace{7\baselineskip}
\subsection{Representación de trenzas sobre \(\mathbb Z/6\mathbb Z\) y caracteres topológicos}
Sea

\[
 w_6:B_n\longrightarrow\mathbb Z/6,
\qquad
 \chi_k(b)=e^{2\pi ikw(b)/6}.
\]
Las representaciones

\[
 \rho_q(\sigma_i)=qP_i
\]
descienden al grupo simétrico únicamente cuando \(q^2=1\). Los otros cuatro valores son genuinamente trenzados. Su existencia matemática no implica cuatro partículas: una realización física requiere Hamiltoniano, brecha espectral, operadores de trenza, estabilidad y entrelazador con una multisección.

\Needspace{7\baselineskip}
\subsection{Imposibilidad de un taquión persistente}
La recta de acción es positiva,

\[
 \hbar_{\rm int}>0,
\]
el carácter satisface

\[
 \chi_{\rm full}>0,
\]
y el operador beta se obtiene por conjugación positiva:

\[
 \widehat M^\beta
 =R_{\rm act}^{\widehat K/2}
 \widehat M^{(0)}
 R_{\rm act}^{\widehat K/2}\ge0.
\]
Por tanto,

\[
 \operatorname{Spec}(\widehat M^\beta)\subset[0,\infty).
\]
Una sección persistente estable no puede generar

\[
 m^2<0.
\]
Si la linealización de un estado produce curvatura negativa o frecuencia imaginaria, ese estado es una inestabilidad que abandona la sección, no una partícula superlumínica. Si la prolongación termina, es una sección virtual con obstrucción terminal. Así, “taquión” se descompone en dos objetos HMT bien tipados ---inestabilidad o virtualidad--- y no constituye una especie persistente.

\Needspace{7\baselineskip}
\subsection{Gravedad sin partícula elemental obligatoria de espín \(2\)}
La realización gravitatoria sigue la cadena

\[
 \text{holonomía de ruta}
 \longrightarrow(e,\omega)
 \longrightarrow(T,R)
 \longrightarrow
 S_{\rm PCH},
\]
donde \(e\) es la coestructura, \(\omega\) la conexión,

\[
 T=de+\omega\wedge e,\qquad
 R=d\omega+\omega\wedge\omega,
\]
y \(S_{\rm PCH}\) es la acción de Palatini--Cartan--Holst. La gravedad aparece como respuesta global de torsión y curvatura de la red persistente. La vibración colectiva del universo modelo puede poseer una componente tensorial de espín dos, pero HMT no necesita postular una partícula elemental adicional para transportar la gravedad.

La constante gravitatoria interna se tipa mediante

\[
 G_{\rm HMT}
 =\frac{c_{\rm int}^3}{\hbar_{\rm int}}
 \bigl(\delta_\theta\Lambda_G\bigr)^2,
\]
donde \(\delta_\theta\) es la torsión angular elemental y \(\Lambda_G\) la longitud espectral seleccionada por la holonomía gravitatoria. Esta ecuación convierte torsión y escala de ruta en acoplamiento gravitatorio sin usar una masa como normalización. La vibración total se representa por el Hessiano

\[
 \mathcal H_{\rm univ}
 =\operatorname{Hess}_{(e,\omega)}S_{\rm PCH}
\]
sobre el estado global. Sus modos son excitaciones colectivas del universo modelo; una componente tensorial no es automáticamente una partícula separada.

Una excitación que se denomine “gravitón” debe demostrar tres cosas:

\begin{enumerate}
\item un subespacio tensorial persistente;
\item un polo aislado del operador linealizado;
\item un acoplamiento físico que no duplique la torsión ya presente.
\end{enumerate}
Sin esas tres condiciones, el objeto correcto es el modo colectivo de la geometría, no una nueva fila de masa.

\Needspace{7\baselineskip}
\subsection{Barbero–Immirzi y el enlace hexada–octada}
Con

\[
 A=1000\alpha_{\rm HMT},
\qquad
 q_\pm=\exp\!\left[-\frac{\pi}{180}(A\pm C^*)\right],
\]
el funcional hexada--octada usa dos incidencias distintas de las torres temporales. La combinación

\[
 \varepsilon_{\rm BI}
 =12\Delta S_{90}-\Delta S_{120}
\]
y el término angular producen

\[
 \gamma_{\rm BI}
 =\frac{\pi}{180}AC^*
 +\varepsilon_{\rm BI}
 =0.2740076726944593\ldots.
\]
El peso \(12:-1\) procede de la lectura de los doce sectores dodecafásicos
y de la única costura orientada de retorno. La ecuación
\(4a-3b=45\) verifica posteriormente que el par \((12,1)\) es la solución
entera positiva mínima. Este objeto no debe confundirse con el extractor dodecafásico de masa ni con el residuo angular \(90\)--\(120\); comparte el operador armónico, pero posee otro dominio y otra función.

\Needspace{7\baselineskip}
\subsection{Extensiones del atlas hacia sectores oscuros: dominios, lectores y estatuto físico}
Un sector oscuro HMT es una multisección persistente con

\[
 P_{\rm vis}P_{\rm dark}=0
\]
para los acoplamientos electromagnético y cromático, pero con posible incidencia gravitatoria o de frontera:

\[
 P_{\rm grav}P_{\rm dark}\ne0.
\]
Su “oscuridad” es una propiedad del núcleo de acoplamiento, no una etiqueta de masa. La ley predice primero un operador y su espectro; sólo un canal físico determina si el modo es estable, metaestable o resonante. Esta definición permite buscar multisecciones invisibles sin inventar una masa mediante el valor observado de una densidad cosmológica.

\Needspace{7\baselineskip}
\subsection{Sección neutrínica estéril: compatibilidad de fibra, mezcla y ausencia de carga de calibre}
Un neutrino estéril no es la profundidad \(G_4\). Es un autovector \(n_s\in\mathcal H_\nu\) que satisface

\[
 P_{\rm weak}n_s=0,\qquad
 P_{\rm mass}n_s=n_s.
\]
Puede mezclarse con los tres modos activos mediante un núcleo neutral ampliado, pero su existencia requiere

\[
 \dim\ker(P_{\rm weak}|_{\mathcal H_\nu})>0
\]
y un autovalor de masa en ese núcleo. Esta definición produce un criterio de búsqueda interno: diagonalizar primero el bloque neutral y comprobar después su acoplamiento débil.

\Needspace{7\baselineskip}
\subsection{Secciones axiónica y pseudoescalares: carácter topológico y acoplamiento de borde}
Un candidato axiónico es una sección pseudoescalar cuya orientación cambia bajo CP,

\[
 CP\,a\,CP^{-1}=-a,
\]
y cuyo desplazamiento compensa una fase topológica del sector de color. En HMT debe corresponder a un modo de torsión/holonomía impar y a un proyector que entrelace ese modo con la densidad topológica de calibre. Su masa se obtiene del segundo derivado del potencial efectivo,

\[
 m_a^2=
 \left.\frac{\partial^2V_{\rm eff}}{\partial a^2}\right|_{a=a_0},
\]
una vez que \(V_{\rm eff}\) haya sido construido desde el registro genealógico de vacancias y frontera. La mera existencia de un canal \(C^*\mapsto-C^*\) no basta para afirmar una partícula axiónica.

\Needspace{7\baselineskip}
\subsection{Fotón oscuro y vector masivo oculto}
Un fotón oscuro requiere una segunda sección de calibre \(A'\) y un acoplamiento de mezcla

\[
 \mathcal L_{\rm mix}
 =-\frac{\varepsilon_{\rm kin}}2F_{\mu\nu}F'^{\mu\nu}.
\]
En términos de fibras, \(\varepsilon_{\rm kin}\) es un entrelazador entre dos subespacios nulos. Si el segundo sector adquiere un brecha espectral Stückelberg o de frontera,

\[
 \widehat M_{A'}^2\ne0,
\]
aparece un vector masivo oscuro. HMT predice que el parámetro de mezcla debe proceder de una incidencia de memoria común; no puede fijarse con la masa que se pretende explicar.

\Needspace{7\baselineskip}
\subsection{WIMP y estado oscuro estable}
Un WIMP es una sección oscura que además cumple:

\[
 [P_{\rm weak},P_X]\ne0,\qquad
 [P_{\rm em},P_X]=[P_{\rm color},P_X]=0,
\]
y posee un modo estable

\[
 |\lambda_X|=1.
\]
La estabilidad puede estar protegida por una paridad del registro genealógico. Su masa es un autovalor del bloque oscuro acoplado débilmente. “WIMP” no constituye una nueva familia matemática: es una condición conjunta de acoplamiento y persistencia sobre una multisección.

\Needspace{7\baselineskip}
\subsection{Sección monopolar: carga magnética, condición de Dirac y compatibilidad topológica}
Un monopolo se define por una clase topológica no trivial del fibrado de calibre:

\[
 \frac1{2\pi}\int_{S^2}F=n\ne0.
\]
HMT debe realizar este entero como holonomía neta de una familia cerrada de rutas. Su masa, si el modo es dinámico y estable, pertenece al espectro del operador de defecto topológico. Una vacancia puntual o un acarreo no nulo no son por sí solos un monopolo; debe demostrarse el descenso a la clase integral.

\Needspace{7\baselineskip}
\subsection{Extensiones supersimétricas del atlas: emparejamiento de representaciones y caracteres}
Una supersimetría HMT requeriría un operador impar \(Q\) que intercambie las completaciones exterior y simétrica:

\[
 Q:\mathcal F_-\leftrightarrow\mathcal F_+,
\qquad
 Q^2=H,
\]
y que entrelace la ley de masas,

\[
 Q\widehat M_-=\widehat M_+Q.
\]
Sólo entonces los espectros bosónico y fermiónico quedarían emparejados. La existencia separada de CAR y CCR no demuestra supersimetría ni predice por sí misma selectrones, neutralinos o gluinos. Cada compañero exige un entrelazador no nulo y una regla de ruptura que determine su separación de masa.

\Needspace{7\baselineskip}
\subsection{Tetraquarks, pentaquarks, glueballs e híbridos}
La gramática de composición se prolonga sin cambiar de ley:

\[
 \mathfrak C_{12}^{(n)}:
 \Gamma_1\times\cdots\times\Gamma_n
 \times\mathcal T_n\times\partial_n
 \longrightarrow\mathcal L_{\rm comp}.
\]
Un tetraquark usa \(qq\bar q\bar q\); un pentaquark, \(qqqq\bar q\); un glueball, registros de la adjunta de color proyectados al singlete; un híbrido, rutas quark y de calibre en un mismo árbol. La clasificación física depende del proyector de color, espín, paridad y árbol de enlace, no sólo de la lista de constituyentes.

\Needspace{7\baselineskip}
\subsection{Secciones de dilatón y radión: modos escalares de volumen y compactificación}
El dilatón es un modo de la sección de escala; el radión, un modo que modifica una longitud o volumen interno. En la teoría de masa aparecen como fluctuaciones de la carta y de la memoria geométrica:

\[
 t_0\mapsto e^{\phi_d}t_0,\qquad
 \mathcal V_{\rm int}\mapsto e^{\phi_r}\mathcal V_{\rm int}.
\]
Para convertirse en partículas deben poseer término cinético, potencial, proyector persistente y polo. Mientras no los posean son grados colectivos de la carta dimensional, no especies escalares adicionales.

\Needspace{7\baselineskip}
\subsection{Exclusión de fantasmas físicos}
La medida espectral de un estado físico es positiva:

\[
 \mu_{P,a}^{\rm mass}(B)\ge0.
\]
Una dirección de norma negativa no pertenece al espacio de estados preparado. Si una parametrización introduce un “fantasma”, la realización física debe cocientar la redundancia de calibre o rechazar el acoplamiento. Esta exclusión es independiente de la prohibición taquiónica: un fantasma viola positividad de norma; un taquión persistente violaría positividad del operador de masa.

\Needspace{7\baselineskip}
\section{Inventario universal de masas y anchuras}
\Needspace{7\baselineskip}
\subsection{Correspondencia posterior entre las realizaciones HMT–MD y el censo físico externo}
El inventario de contraste contiene

\[
 1170\ \text{identidades},\qquad
 438\ \text{grupos de identidad},
\]
\[
 471\ \text{observables de masa},\qquad
 384\ \text{observables de anchura}.
\]
Su partición es

\Needspace{7\baselineskip}
\begingroup
\footnotesize
\setlength{\tabcolsep}{2.1pt}
\begin{longtable}{L{5.04cm}L{5.04cm}L{5.04cm}}
\toprule
\rowcolor{HMTNavy}\HMTtablehead{Dominio} & \HMTtablehead{Masas} & \HMTtablehead{Anchuras} \\
\midrule
\endfirsthead
\multicolumn{3}{l}{\sffamily\scriptsize\color{HMTGray}Continuación de la tabla}\\[-1pt]
\toprule
\rowcolor{HMTNavy}\HMTtablehead{Dominio} & \HMTtablehead{Masas} & \HMTtablehead{Anchuras} \\
\midrule
\endhead
\midrule
\endfoot
\bottomrule
\endlastfoot
partículas generales y leptónicas & 53 & 9 \\
\rowcolor{HMTLight}quarks & 8 & 1 \\
mesones & 243 & 225 \\
\rowcolor{HMTLight}bariones & 166 & 149 \\
registro gravitatorio de contraste & 1 & 0 \\
\rowcolor{HMTLight}\textbf{Total} & \textbf{471} & \textbf{384} \\
\end{longtable}
\endgroup
Las 471 masas no son 471 especies elementales: incluyen estados excitados, resonancias, distintas cargas y múltiples realizaciones compuestas. Las 384 anchuras son observables temporales y no deben ser leídas por el carácter de masa.

\Needspace{7\baselineskip}
\subsection{Identidad interna \(\mathcal I_X\) de una realización material y criterio de individualización frente a observaciones múltiples}
Cada estado comparado debe poseer una identidad interna

\[
 \mathcal I_X=
 (\mathcal C_X,\mathcal R_X,\Gamma_X,\mathcal L_X,
 \sigma_X,\eta_X,P_X,\mathfrak M_X,\mathcal T_X),
\]
donde:

\begin{itemize}
\item \(\mathcal C_X\) es el tipo físico de realización determinado por objeto, representación y codominio;
\item \(\mathcal R_X\) es la representación física;
\item \(\Gamma_X\) es una ruta, órbita o composición orientada;
\item \(\mathcal L_X\) es el registro genealógico;
\item \(\sigma_X\) es la firma completa;
\item \(\eta_X\) es la torre;
\item \(P_X\) es el proyector del acoplamiento;
\item \(\mathfrak M_X\) es el canal de medición;
\item \(\mathcal T_X\) es el árbol de composición cuando procede.
\end{itemize}
Dos observaciones externas pertenecen al mismo estado HMT cuando comparten esta identidad y difieren sólo por experimento, carta o estimador. No se crean dos partículas porque existan dos mediciones.

\Needspace{7\baselineskip}
\subsection{Morfismo de comparación entre invariantes internos y observables metrológicos}
La comparación requiere el diagrama

\[
 \mathcal L_M^{\rm HMT}
 \xrightarrow{\;\varsigma_{\mathcal S,\mu}\;}
 \mathbb R_{>0}u_M^{\mathcal S,\mu}
 \xleftarrow{\;\text{medición}\;}
 \mathcal O_X^{\rm ext}.
\]
Una fila científicamente completa contiene:

\[
 (m_X^{\rm HMT},m_X^{\rm ext},
 u,\mathcal S,\mu,\Sigma_{\rm exp},\Sigma_{\rm HMT},
 \text{convención de polo}).
\]
Para quarks son obligatorios esquema y escala; para resonancias, la convención de polo; para límites, el nivel de confianza; para observables correlacionados, la matriz de covarianza.

\Needspace{7\baselineskip}
\subsection{Evaluación explícita de \(17\) realizaciones representativas del atlas}
Las siguientes evaluaciones conservan la separación entre ley y dato de firma:

\Needspace{7\baselineskip}
\begingroup
\footnotesize
\setlength{\tabcolsep}{2.1pt}
\begin{longtable}{L{4.68cm}L{10.61cm}}
\toprule
\rowcolor{HMTNavy}\HMTtablehead{Estado} & \HMTtablehead{Valor del carácter en su carta} \\
\midrule
\endfirsthead
\multicolumn{2}{l}{\sffamily\scriptsize\color{HMTGray}Continuación de la tabla}\\[-1pt]
\toprule
\rowcolor{HMTNavy}\HMTtablehead{Estado} & \HMTtablehead{Valor del carácter en su carta} \\
\midrule
\endhead
\midrule
\endfoot
\bottomrule
\endlastfoot
\(\mu\) & \(105.658360\allowbreak{}294435\allowbreak{}829\ldots\ {\rm MeV}\) \\
\rowcolor{HMTLight}\(\tau\) & \(1776.860917\allowbreak{}631432\allowbreak{}3\ldots\ {\rm MeV}\) \\
\(u\) & \(2.165924\allowbreak{}536173\allowbreak{}907\ldots\ {\rm MeV}\) \\
\rowcolor{HMTLight}\(d\) & \(4.679550\allowbreak{}162948\allowbreak{}874\ldots\ {\rm MeV}\) \\
\(c\) & \(1.270500\allowbreak{}838641\allowbreak{}911\ldots\ {\rm GeV}\) \\
\rowcolor{HMTLight}\(s\) & \(92.940093\allowbreak{}907845\allowbreak{}64\ldots\ {\rm MeV}\) \\
\(t\) & \(172.597479\allowbreak{}544019\allowbreak{}1\ldots\ {\rm GeV}\) \\
\rowcolor{HMTLight}\(b\) & \(4.190119\allowbreak{}763299\allowbreak{}790\ldots\ {\rm GeV}\) \\
\(W\) & \(80.378964\allowbreak{}909704\allowbreak{}55\ldots\ {\rm GeV}\) \\
\rowcolor{HMTLight}\(Z\) & \(91.187533\allowbreak{}444313\allowbreak{}41\ldots\ {\rm GeV}\) \\
\(H\) & \(125.292466\allowbreak{}042536\allowbreak{}1\ldots\ {\rm GeV}\) \\
\rowcolor{HMTLight}\(\pi^0\) & \(134.976724\allowbreak{}327872\allowbreak{}1\ldots\ {\rm MeV}\) \\
\(\pi^\pm\) & \(139.570485\allowbreak{}171572\allowbreak{}2\ldots\ {\rm MeV}\) \\
\rowcolor{HMTLight}\(K^\pm\) & \(493.677085\allowbreak{}649530\allowbreak{}6\ldots\ {\rm MeV}\) \\
\(K^0\) & \(497.611186\allowbreak{}497750\allowbreak{}5\ldots\ {\rm MeV}\) \\
\rowcolor{HMTLight}\(p\) & \(938.271751\allowbreak{}546984\allowbreak{}9\ldots\ {\rm MeV}\) \\
\(n\) & \(939.565810\allowbreak{}472507\allowbreak{}0\ldots\ {\rm MeV}\) \\
\end{longtable}
\endgroup
El electrón beta no pertenece a esta tabla porque su ruta central, sus coeficientes basales y su lector completo están determinados por la propia construcción.

\Needspace{7\baselineskip}
